\documentclass[12pt, titlepage]{article}
\usepackage{amsmath}
\usepackage{amsfonts}
\usepackage{amssymb}
\usepackage{fancyhdr}
\usepackage[bidi=basic]{babel}
\babelprovide[main,import]{hebrew}
\babelfont[hebrew]{rm}{Hadasim CLM}
\babelfont[hebrew]{sf}{Miriam Mono CLM}
\usepackage{titlesec}
\title{מושגי יסוד במתמטיקה 01040002 מועד א' חורף 2025}
\author{יונתן אבידור - 214269565}
\titleformat{\section}{\Large\sffamily\bfseries}{\thesection}{1em}{}
\begin{document}
\maketitle
\pagestyle{fancy}
\fancyhead{}
\fancyfoot[L]{\text{214269565}}
\part*{שאלה 1}
\section*{סעיף א'}
$\left[a\right]_R$, המכונה גם מחלקת השקילות של $a$ תחת $R$, היא קבוצה המכילה את כל האיברים $b\in A$ כך ש $\left(a,b\right) \in R$, כלומר כל האיברים ב$A$ שעומדים ביחס $R$ עם $a$\\
$A\setminus R$, המכונה גם קבוצת המנה של $R$ היא קבוצת כל מחלקות השקילות של $A$, כלומר
\[
    R=\bigl\{\left[a\right]_R~|~a\in A \bigr\}
\]
\section*{סעיף ב'}
\subsection*{i.}
נתון כי $S_1 = S_2$ ו-$S_2$ הוא חתך של $R_2$, לכן $S_1$ הוא גם חתך של $R_2$.\\
נטען כי $\left|S_1\right|=\left|A\setminus R_1\right|$ וכי $\left|S_1\right|=\left|A\setminus R_2\right|$
\subsubsection*{טענת עזר}
תהא קבוצה $A$, יהי $R$ יחס שקילות על $A$, יהי $S$ חתך של $R$.\\
אזי
\[
 \left|S\right|=\left|A\setminus R\right|
\]
\subsubsection*{הוכחת הטענה}
נבנה פונקציה חד-חד ערכית ועל מהחתך לקבוצת המנה
\begin{gather*}
    f:S\to A\setminus R\\
    \forall s\in S:    f\left(s\right)=\left[s\right]_R
\end{gather*}
נוכיח כי היא חד-חד ערכית ועל\\
חד-חד ערכית\\
יהיו $s_1,s_2\in S$ כך ש-$f\left(s_1\right) = f\left(s_2\right)$, נרצה להראות ש-$s_1=s_2$
\[
    f\left(s_1\right)=f\left(s_2\right)\implies \left[s_1\right]_R = \left[s_2\right]_R
\]
מהגדרת חתך, לכל מחלקת שקילות יש נציג אחד בלבד\\
לכן אם שני איברים בחתך הם מאותה מחלקת שקילות, הם אותו האיבר.\\
מכאן $s_1 = s_2$\\
על:\\
תהא $X \in A\setminus R$, נרצה להראות כי $\exists s\in S: f\left(s\right)=X$\\
מהגדרת חתך, \pmb{לכל} מחלקת שקילות קיים נציג בחתך, מכיוון שהפונקציה ממפה את איברי החתך למחלקת השקילות שלהם, לכל מחלקת השקילות בהכרח יהיה איבר בחתך שנמצא באותה מחלקת שקילות, ולכן הוא מקור אליה תחת $f$
\bigskip\\
נחזור כעת להוכחה עצמה.\\
על פי טענת העזר, עוצמת קבוצת המנה של יחס שווה לעוצמת חתך שלה\\
מהנתון, $S_1$ חתך של $R_1$ ולכן $\left|S_1\right|=\left|A\setminus R_1\right|$, בנוסף, הראינו כי $S_1$ חתך של $R_2$ ולכן, מטענת העזר $\left|S_1\right|=\left|A\setminus R_2\right|$
\[
    \left|S_1\right|=\left|A\setminus R_1\right|\land\left|S_1\right|=\left|A\setminus R_2\right|\implies\left|A\setminus R_1\right|=\left|A\setminus R_2\right|
\]
כנדרש\\
\blacksquare
\subsubsection*{ii.}
לא נכון, נביא דוגמה נגדית
\begin{gather*}
    A=\{1,2,3\}\\
    R_1 = \left\{\left(1,1\right),\left(2,2\right),\left(3,3\right),\left(1,2\right),\left(2,1\right)\right\}\\
    R_2 = \left\{\left(1,1\right),\left(2,2\right),\left(3,3\right),\left(3,2\right),\left(2,3\right)\right\}\\
    S_1= \left\{1,3\right\}\\
    S_2= \left\{1,3\right\}
\end{gather*}
$S_1=S_2$ כי בשניהם יש אך ורק את האיברים $1,3$, אבל $R_1\neq R_2$זה מכיוון ש $\left(1,2\right)\in R_1$  אבל $\left(1,2\right) \notin R_2$, קיים ב-$R_1$ איבר שלא נמצא ב-$R_2$ ולכן $R_1 \not\subseteq R_2$ ולכן $R_1 = R_2$
\part*{שאלה 2}
\section*{סעיף א'}
$f$ אינה חח"ע\\
נראה באמצעות שתי קבוצות $A,B\in\mathcal{P}\left(\mathbb{N}\right)$ כך ש$A\neq B$ אבל $f(A)=f(B)$
\begin{gather*}
    A=\{1\}\\
    B=\emptyset\\
    1\notin \emptyset \implies A\not\subseteq B\implies A \neq B\\
    f(A)=A\setminus\{1\}=\{1\}\setminus\{1\}=\emptyset\\
    f(B)=B\setminus\{1\}=\emptyset\setminus\{1\}=\emptyset
\end{gather*}
$f(A)=f(B)$ אבל $A\neq B$ ולכן $f$ לא חח"ע
\section*{סעיף ב'}
נרצה שסכום או הפרש של שני מספרים יהיה כפולה של $10$\\
על מנת שהסכום או ההפרש של שני מספרים יהיה 10, יש 6 אפשרויות לפי השארית החלוקה של המספרים עם $10$
\begin{gather*}
    \boxed{\left(0,10\right)~\left(10,0\right)}~~\boxed{\left(1,9\right)~\left(9,1\right)}~~\boxed{\left(2,8\right)~\left(8,2\right)}\\
    \boxed{\left(3,7\right)~\left(7,3\right)}~~\boxed{\left(4,6\right)~\left(6,4\right)}~~\boxed{\left(5,5\right)}
\end{gather*}
על פי עיקרון שובך היונים, יש לנו $6$ "תאים" ו$7$ "יונים", אזי בהכרח באחד מהם יהיו שני "יונים"
\part*{שאלה 3}
\section*{סעיף א'}
$\left|X\right|<\left|Y\right|$ אומר ש $|X|\leq|Y|$ וגם $|X|\neq|Y|$, כלומר שקיימת פונקציה חח"ע מ-$X$ ל-$Y$ וגם שלא קיימת פונקציה חח"ע ועל מ-$X$ ל$Y$
\end{document}
